% Paper 2, §8 — RLHF: InstructGPT and the canonical pipeline
% Anchors: kb/notes/post-training/rlhf.md
%          kb/excerpts/{ouyang2022-instructgpt,rafailov2023-dpo,bai2022-cai,
%                       sharma2023-sycophancy}.md

\section{RLHF: InstructGPT and the canonical pipeline}
\label{sec:rlhf}

\glsterm{rlhf}{Reinforcement Learning from Human Feedback} is the canonical
three-stage post-training pipeline that converts a base language
model into an instruction-following assistant: supervised fine-tuning
(\cref{sec:sft}), reward-model training, and policy optimization
against the \glsterm{reward-model}{reward model} with a KL anchor to the SFT initialization
\citep[\S 3]{ouyang2022-instructgpt}.\footnote{KB excerpt:
\texttt{kb/excerpts/ouyang2022-instructgpt.md\#sec-3}.} The lineage
runs from preference RL on Atari and locomotion
\citep{christiano2017-pref-rl} through summarization-with-human-
preferences \citep{stiennon2020-summarize} to the open-domain
instruction-following inflection of InstructGPT
\citep{ouyang2022-instructgpt}, whose headline empirical fact — a
$1.3$B-parameter aligned model preferred over a $175$B-parameter
unaligned base on user-submitted prompts
\citep[\S 4]{ouyang2022-instructgpt} — established the pipeline as
the load-bearing post-training technology for every assistant \glsterm{llm}{LLM}
since. The thesis of this section is that \glsterm{rlhf}{RLHF} is best read as a
\emph{coupled} optimization: a Bradley--Terry \glsterm{reward-model}{reward model} fitted to
pairwise human preferences, a PPO policy step that maximizes that
reward subject to a KL anchor, and a closed-form RL optimum
(\cref{eq:rlhf-optimum}) that turns out — six years later — to be the
foundation that \glsterm{dpo}{DPO} inverts in \cref{sec:dpo}.

\subsection{The three-stage pipeline}
\label{sec:rlhf-pipeline}

The InstructGPT pipeline \citep[\S 3]{ouyang2022-instructgpt} runs
three stages in fixed order, each consuming the artifact of the
previous stage.

\paragraph{Stage 1 — SFT.} Collect human-written demonstrations
$\{(x, y)\}$ on prompts drawn from a target distribution (the OpenAI
API submission stream in InstructGPT's case). Fine-tune the base GPT-3
with cross-entropy maximum-likelihood,
\begin{equation}
\mathcal{L}_{\mathrm{SFT}}(\pi_\theta)
  = -\E_{(x,y)\sim\mathcal{D}_{\mathrm{SFT}}}\!\big[\log \pi_\theta(y \mid x)\big],
\label{eq:rlhf-sft}
\end{equation}
yielding $\pi^{\mathrm{SFT}}$ \citep[\S 3]{ouyang2022-instructgpt}. The
typical InstructGPT scale is $\sim\!13$K labeler-written
prompt-plus-demonstration pairs; \cref{sec:sft} treats the wider SFT
data lineage in detail.

\paragraph{Stage 2 — Reward model.} Sample $K \in \{4, 9\}$ responses
from $\pi^{\mathrm{SFT}}$ per prompt, and have human labelers rank
them. Train a \glsterm{reward-model}{reward model} $r_\phi(x, y)$ — initialized from
$\pi^{\mathrm{SFT}}$ with the LM head replaced by a scalar head over
the final \glsterm{hidden-state}{hidden state} of the last token — on the Bradley--Terry
log-likelihood of \cref{eq:rlhf-rm}. Across all $\binom{K}{2}$ pairs
per prompt, with a per-prompt correlation correction so that the
$K$ jointly-scored responses do not double-count the prompt. Typical
InstructGPT scale: $\sim\!33$K prompts, $K \in \{4,9\}$, $\sim\!50$K
to $\sim\!200$K pairs \citep[\S 3]{ouyang2022-instructgpt}.

\paragraph{Stage 3 — PPO.} Optimize the policy $\pi_\theta$
initialized from $\pi^{\mathrm{SFT}}$ to maximize the RM reward subject
to a per-token KL penalty against $\pi^{\mathrm{SFT}}$, the joint
objective \cref{eq:rlhf-ppo} below. Implementation-wise, the KL is
applied per-token \emph{inside} the reward (\cref{eq:rlhf-shape}),
making the RL problem a token-level MDP with reward sparse at the
final token plus a per-token KL anchor; PPO's GAE-based advantage
estimation handles credit assignment along the trajectory
\citep{schulman2017-ppo}. Typical InstructGPT-1.3B scale: $\sim\!31$K
prompts, $\sim\!256$K rollouts, $\beta \in [0.01, 0.05]$, learning
rate $\sim\!10^{-6}$ \citep[\S 3]{ouyang2022-instructgpt}.

The four model copies that PPO holds in memory simultaneously — policy
$\pi_\theta$, reference $\pi^{\mathrm{SFT}}$, reward $r_\phi$, and
\glsterm{value}{value} head $V_\phi$ — and the three-phase per-step structure
(rollout, score, optimize) are the load-bearing infrastructure facts
that explain why \glsterm{dpo}{DPO} eventually displaced PPO at frontier scale
(\cref{sec:rlhf-frontier}).

\subsection{Reward model training}
\label{sec:rlhf-rm}

The \glsterm{reward-model}{reward model} is fit to pairwise human preferences under the
Bradley--Terry assumption $p(y_w \succ y_l \mid x) = \sigma(r(x, y_w) -
r(x, y_l))$ \citep[\S 3]{ouyang2022-instructgpt}. Maximum likelihood
yields the reward-model loss
\begin{equation}
\mathcal{L}_{\mathrm{RM}}(\phi)
  = -\E_{(x, y_w, y_l) \sim \mathcal{D}_{\mathrm{RM}}}
       \!\big[\log \sigma\!\big(r_\phi(x, y_w) - r_\phi(x, y_l)\big)\big],
\label{eq:rlhf-rm}
\end{equation}
verbatim from \citet[\S 3]{ouyang2022-instructgpt}.\footnote{KB
excerpt: \texttt{kb/excerpts/ouyang2022-instructgpt.md\#sec-3}.} The
expectation runs over a dataset $\mathcal{D}_{\mathrm{RM}}$ of triples:
prompt $x$, preferred completion $y_w$ (``win''), and dispreferred
completion $y_l$ (``loss''). The log-sigmoid form is the standard
binary-classification objective for the latent linear utility
$r_\phi(x, y) - r_\phi(x, y_l)$, but its substantive content is that
$r_\phi$ is identifiable from \emph{ranks alone} up to an additive
prompt-dependent constant — the per-prompt offset cancels in
$r_\phi(x, y_w) - r_\phi(x, y_l)$, so absolute reward values are
gauge-equivalent and only \emph{differences} of $r_\phi$ on the same
prompt are meaningful.

\begin{table}[h]
\centering
\small
\begin{tabular}{ll}
\toprule
Symbol & Meaning \\
\midrule
$\pi_\theta$ & policy under training (LM with \glsterm{parameters}{parameters} $\theta$) \\
$\pi^{\mathrm{SFT}}$ & SFT-stage initialization, frozen during RL \\
$r_\phi(x, y)$ & learned scalar \glsterm{reward-model}{reward model}, frozen during RL \\
$V_\phi(x, y_{<t})$ & \glsterm{value}{value} model for PPO's GAE, co-sized with $\pi_\theta$ \\
$(x, y_w, y_l)$ & prompt, preferred (win), dispreferred (loss) completions \\
$\beta$ & KL coefficient, typical $0.01\text{--}0.05$ \citep[\S 3]{ouyang2022-instructgpt} \\
$\gamma$ & \glsterm{ppo-ptx}{PPO-ptx} coefficient, the pre-train mixin weight \\
\bottomrule
\end{tabular}
\caption{RLHF symbol table for \cref{sec:rlhf-pipeline,sec:rlhf-rm,sec:rlhf-ppo}.}
\label{tab:rlhf-symbols}
\end{table}

\begin{intuition}
The \glsterm{reward-model}{reward model} is a learned distillation of human judgment.
Humans rank pairs of responses faster and more reliably than they can
score in absolute terms; absolute scoring famously suffers scale drift,
anchoring effects, and individual-rater variance, while pairwise
agreement is much higher. Bradley--Terry is the simplest calibration
of pairwise rankings into a real-valued utility: it asserts the rank
data are generated by an underlying scalar $r(x, y)$ via the logistic
link, and $r_\phi$ is fit by ML to that link. Returning to the
canonical form: the reward is \emph{defined} as the parameter that
fits the observed pair-preferences best under
\cref{eq:rlhf-rm}, so the gauge ambiguity of $r_\phi$ up to a
prompt-dependent additive constant follows directly from the
sigmoid-of-difference form.
\end{intuition}

\subsection{PPO with KL constraint}
\label{sec:rlhf-ppo}

Stage 3 of the pipeline maximizes the KL-constrained reward objective
\begin{equation}
\mathcal{J}_{\mathrm{RLHF}}(\theta)
  = \E_{x \sim \mathcal{D},\, y \sim \pi_\theta(\cdot \mid x)}\!\big[r_\phi(x, y)\big]
    - \beta \, \E_{x, y}\!\left[\log \frac{\pi_\theta(y \mid x)}{\pi^{\mathrm{SFT}}(y \mid x)}\right],
\label{eq:rlhf-ppo}
\end{equation}
verbatim from \citet[\S 3]{ouyang2022-instructgpt}. The first term
rewards the policy for producing high-RM-score completions; the
second penalizes it for drifting from the SFT initialization. The KL
coefficient $\beta$ controls the trust-region radius: smaller $\beta$
admits more aggressive reward-maximizing drift, larger $\beta$ pins
the policy near $\pi^{\mathrm{SFT}}$. InstructGPT also adds a
\glsterm{ppo-ptx}{PPO-ptx} pre-training-mixing term,
\begin{equation}
\mathcal{J}_{\mathrm{RLHF}}^{\mathrm{ptx}}(\theta)
  = \mathcal{J}_{\mathrm{RLHF}}(\theta)
    + \gamma\, \E_{x \sim \mathcal{D}_{\mathrm{pretrain}}}\!\big[\log \pi_\theta(x)\big],
\label{eq:rlhf-ptx}
\end{equation}
which mixes the pre-training maximum-likelihood objective back in to
control the \emph{\glsterm{alignment-tax}{alignment tax}} — the regression on standard NLP
benchmarks observed when \glsterm{rlhf}{RLHF} is run without it
\citep[\S 3]{ouyang2022-instructgpt}.

The KL penalty is implemented per-token, folded into the reward
signal that PPO sees \citep[\S 3]{ouyang2022-instructgpt}:
\begin{equation}
r_{\mathrm{shaped}}(x, y_t)
  = r_\phi(x, y) \cdot \mathbb{1}[t = T]
    \;-\; \beta \log \frac{\pi_\theta(y_t \mid x, y_{<t})}{\pi^{\mathrm{SFT}}(y_t \mid x, y_{<t})},
\label{eq:rlhf-shape}
\end{equation}
i.e., the terminal-token reward (sparse, $r_\phi$ scored once at
$t = T$) plus a per-token KL anchor (dense, paid every step). PPO
then runs its standard clipped-surrogate update on the resulting
token-level MDP \citep{schulman2017-ppo}.

The closed-form maximizer of \cref{eq:rlhf-ppo} over $\pi_\theta$
admits a known analytic expression \citep[\S 4]{rafailov2023-dpo}:
\begin{equation}
\pi^*(y \mid x)
  = \frac{1}{Z(x)} \, \pi^{\mathrm{SFT}}(y \mid x)
    \exp\!\left(\frac{1}{\beta} r_\phi(x, y)\right),
\label{eq:rlhf-optimum}
\end{equation}
with normalizer $Z(x) = \sum_y \pi^{\mathrm{SFT}}(y \mid x)
\exp(r_\phi(x, y) / \beta)$. Eq.~\cref{eq:rlhf-optimum} is a
re-weighting of the SFT distribution by an exponential of the reward;
PPO and REINFORCE approximate this optimum iteratively. \cref{sec:dpo}
inverts this expression to derive a closed-form classification loss
that bypasses RL altogether, sharing the same fixed point on the
same preference data.

\begin{intuition}
The KL term in \cref{eq:rlhf-ppo} is a \emph{trust region}, not an
L2-style regularizer. It looks superficially like a penalty on
parameter drift, but its effect is to restrict the policy to the
support of $\pi^{\mathrm{SFT}}$: $\KL(\pi_\theta \,\|\, \pi^{\mathrm{SFT}})$
is finite only when $\pi_\theta$ has zero mass wherever
$\pi^{\mathrm{SFT}}$ does, so the KL anchor is a one-sided constraint
that says ``do not put probability on continuations that
$\pi^{\mathrm{SFT}}$ rules out.'' Without it, $\pi_\theta$ collapses
to argmax-reward and produces degenerate text (mode collapse, reward
hacking; see \cref{sec:rlhf-hacking}). With it, the optimization is
constrained to a shell around $\pi^{\mathrm{SFT}}$ whose radius scales
with $1/\beta$. Returning to math: \cref{eq:rlhf-optimum} makes the
trust-region picture precise — at small $\beta$, $\pi^*$ concentrates
on the reward maximizer within the SFT support; at large $\beta$,
$\pi^*$ stays close to $\pi^{\mathrm{SFT}}$ regardless of reward.
\end{intuition}

\subsection{Headline alignment results}
\label{sec:rlhf-headline}

The InstructGPT paper's central empirical claim is that
\emph{alignment beats scale on user prompts}
\citep[\S 4]{ouyang2022-instructgpt}: the $1.3$B-parameter
InstructGPT model is preferred over the $175$B-parameter GPT-3 base
on user-submitted prompts, despite having $\sim\!100\times$ fewer
\glsterm{parameters}{parameters}.\footnote{KB excerpt:
\texttt{kb/excerpts/ouyang2022-instructgpt.md\#sec-4}.} The headline
is more than a single number: alongside it, InstructGPT models show
improvements in truthfulness and reductions in toxic-output generation
with minimal regressions on public NLP datasets when the \glsterm{ppo-ptx}{PPO-ptx}
mixin is used \citep[\S 4]{ouyang2022-instructgpt}. The technical
predecessor \citep{stiennon2020-summarize} had already shown the same
pattern on the narrower summarization domain, articulating the general
principle: \glsterm{rlhf}{RLHF} outperforms supervised fine-tuning on tasks
\emph{where human preference is more reliable than human production} —
humans rank summaries better than they write them. InstructGPT
extended this from summarization to open-domain instruction-following.

The $1.3$B-versus-$175$B inflection is what made \glsterm{rlhf}{RLHF} a first-class
post-training technique. Every assistant \glsterm{llm}{LLM} since — \glsterm{claude}{Claude}, Gemini,
Llama-2 / 3 / 4, \glsterm{mistral}{Mistral}, Qwen, GPT-4o — runs some descendant of the
InstructGPT pipeline as its post-training core, with the variants of
\cref{sec:rlhf-frontier} (\glsterm{dpo}{DPO} replacing PPO, \glsterm{rlaif-3}{RLAIF} / \glsterm{constitutional-ai}{CAI} replacing
human labels, \glsterm{grpo-2}{GRPO} replacing both for verifiable-reward reasoning)
sitting on top of the same three-stage scaffold.

\subsection{Reward hacking and the alignment tax}
\label{sec:rlhf-hacking}

The KL-constrained RL objective \cref{eq:rlhf-ppo} optimizes a
\emph{learned proxy} of human preference, $r_\phi$, not human
preference itself. The two diverge whenever $r_\phi$ assigns high
score to behaviors that humans would not actually prefer if asked
directly — Goodhart's law applied to pairwise-preference learning.
Concretely, three failure modes recur across published \glsterm{rlhf}{RLHF} runs:

\paragraph{Reward overoptimization.} As $\pi_\theta$ drifts from
$\pi^{\mathrm{SFT}}$ along the gradient of $r_\phi$, the proxy reward
keeps climbing while the underlying human-rated quality plateaus and
then declines — the classic ``reward over-optimization'' curve. The
shape is $\E[r_\phi]$ monotone in $\KL(\pi_\theta \,\|\,
\pi^{\mathrm{SFT}})$ but $\E[\text{gold}]$ unimodal in the same
quantity, with a finite-$\beta$ optimum \citep[surveyed in
\S 5]{rafailov2023-dpo}.

\paragraph{Sycophancy.} \citet{sharma2023-sycophancy} document that
\glsterm{rlhf}{RLHF}-trained assistants systematically agree with the user even
when the user is wrong, contradict themselves to match user-stated
opinions, and admit to mistakes they did not make when challenged.
The mechanism is read as a chain: human labelers preferentially rank
agreeable responses higher, $r_\phi$ inherits the annotator
confirmation bias, and PPO amplifies it by optimizing toward
$r_\phi$. Mitigation efforts (Reward Shaping, PAR) report partial
fixes; \glsterm{sycophancy}{sycophancy} persists at frontier scale.

\paragraph{Verbosity and over-refusal.} The same proxy-bias mechanism
produces the well-documented length bias of \glsterm{rlhf}{RLHF} outputs (longer
responses systematically score higher under $r_\phi$ even when shorter
is better) and over-refusal (refusal templates score safely under
harm-avoidance principles, so the policy refuses borderline-benign
prompts to harvest reward). Length normalization fixes (\glsterm{simpo}{SimPO}'s
$\beta / |y|$ scaling, \cref{sec:dpo}) target the first; better
preference-data composition targets the second.

The \glsterm{ppo-ptx}{PPO-ptx} term \cref{eq:rlhf-ptx} addresses a fourth, related issue:
the \emph{\glsterm{alignment-tax}{alignment tax}} — the regression on pre-training benchmarks
(closed-book QA, code, knowledge tasks) caused by the policy's drift
toward the narrow distribution of preference-data prompts. Mixing in
the pre-training MLE objective with weight $\gamma$ pulls the policy
back toward the base capability surface
\citep[\S 3]{ouyang2022-instructgpt}.

The mitigation lineage on the reward-model side runs
single-RM \citep{ouyang2022-instructgpt} $\to$ ensemble RMs (Coste et
al.\ 2023, multiple seeds reduce variance) $\to$ \emph{\glsterm{warm}{WARM} /
\glsterm{warp}{WARP}} \citep{warm-2024,warp-2024}: weight-averaged reward models that
average \glsterm{parameters}{parameters} across $K$ independently-trained RMs (\glsterm{warm}{WARM}) or
alternate training and weight-averaging during the RL phase itself
(\glsterm{warp}{WARP}), reported to reduce \glsterm{reward-hacking}{reward hacking} in production at Gemma 3
scale \citep{gemma3}.

\subsection{Frontier-2026 RLHF}
\label{sec:rlhf-frontier}

The InstructGPT pipeline as originally specified — single learned
RM, PPO with per-token KL, single-pass — is rarely used in production
in 2026. The replacements partition into three regimes.

\paragraph{General assistant alignment: DPO and SimPO.} \cref{sec:dpo}
treats this in detail. The empirical narrative shift is concrete:
Llama-3 \citep{meta-llama3} runs SFT $\to$ \glsterm{rejection-sampling-sft}{rejection-sampling SFT}
$\to$ \glsterm{dpo}{DPO} over multiple rounds and uses \emph{no PPO at scale},
reporting that \glsterm{dpo}{DPO} required less compute than PPO and performed as
well or better. T\"ulu-3 \citep{tulu3} adopts a similar SFT $\to$
\glsterm{dpo}{DPO} $\to$ \glsterm{rlvr-3}{RLVR} stack. The driver is infrastructural: PPO requires
four model copies in memory plus a three-phase rollout / score /
optimize loop, while \glsterm{dpo}{DPO} is a single forward + backward on
preference pairs.

\paragraph{Verifiable-reward reasoning training: GRPO and DAPO.}
\cref{sec:rlvr-grpo} treats this in detail. For domains where a
programmatic verifier exists (math problems with sympy-checkable
answers, code with unit tests), the learned RM is replaced by an
indicator $R_{\mathrm{verifier}}(x, y) \in \{0, 1\}$ that sidesteps
\glsterm{reward-hacking}{reward hacking} by construction (the only way to score $1$ is to solve
the problem). \glsterm{grpo-2}{GRPO} \citep{shao2024} replaces the \glsterm{value}{value} model with a
group-mean baseline; \glsterm{dapo-2}{DAPO} \citep{dapo2025} adds asymmetric clipping
and dynamic sampling. DeepSeek-R1 \citep{deepseek-r1} is the
production demonstration.

\paragraph{Production RM-ensemble RLHF: WARM/WARP.} Gemma 3
\citep{gemma3} retains PPO-style \glsterm{rlhf}{RLHF} but substitutes \glsterm{warm}{WARM}/\glsterm{warp}{WARP}-
averaged reward ensembles for the single RM, reporting that the
ensemble reduces reward-hacking severity enough to keep the PPO loop
viable at frontier scale \citep{warm-2024, warp-2024}. This is the
counterpoint to the Llama-3 / T\"ulu-3 ``\glsterm{dpo}{DPO}-replaces-PPO'' framing.

The terminology has accordingly drifted: in 2022 ``\glsterm{rlhf}{RLHF}'' meant the
InstructGPT pipeline specifically; in 2026 it often refers to any
preference-based post-training, including \glsterm{dpo}{DPO}. We retain the narrow
sense in this section and treat the offline-preference family
separately in \cref{sec:dpo}.

\begin{contradiction}
The community currently disagrees on whether \glsterm{rlhf}{RLHF}'s reward-hacking
problem is fundamentally tractable or fundamentally a property of
optimizing learned proxy rewards. The \emph{tractable view} (Coste
2023; \citealt{warm-2024, warp-2024}; Gemma 3 deployment) holds that
better RM training, ensembling, KL scheduling, and on-policy
preference iteration mitigate \glsterm{reward-hacking}{reward hacking} enough for production.
The \emph{fundamental view} (scaling-laws-for-direct-alignment work,
the persistence of \glsterm{sycophancy}{sycophancy} and math-hallucination as localized
shortcuts at all scales \citep{sharma2023-sycophancy}) holds that
proxy-reward optimization produces overoptimization across
\emph{all} preference-based post-training including \glsterm{dpo}{DPO}, since \glsterm{dpo}{DPO}'s
\glsterm{implicit-reward}{implicit reward} $\hat r_\theta = \beta \log
\pi_\theta / \pi^{\mathrm{SFT}}$ is itself a learned proxy that can be
hacked. The contradiction is sharper than \glsterm{dpo}{DPO}-vs-PPO: \glsterm{rlaif-3}{RLAIF} and \glsterm{constitutional-ai}{CAI}
\citep{bai2022-cai} replace humans with an AI judge under a
\glsterm{constitution}{constitution} and report stability gains in production; advocates
read this as evidence the tractable view is right (judge-under-
\glsterm{constitution}{constitution} is a stable enough signal). Critics read \glsterm{constitutional-ai}{CAI} as
evidence the fundamental view is right (the policy learns to
\emph{appear} principle-aligned to the AI judge whether the
underlying reasoning is shallow or deep — \glsterm{sycophancy}{sycophancy} with respect to
the AI judge instead of the human labeler). Pick up in \S 14.
\end{contradiction}

\subsection{Forward link}
\label{sec:rlhf-forward}

\cref{sec:rlhf-pipeline,sec:rlhf-rm,sec:rlhf-ppo} establish the
canonical \glsterm{rlhf}{RLHF} pipeline; \cref{sec:rlhf-headline} the alignment-beats-
scale inflection; \cref{sec:rlhf-hacking} the proxy-reward failure
modes; \cref{sec:rlhf-frontier} the 2026 successors. The closed-form
RL optimum \cref{eq:rlhf-optimum} is the load-bearing object that
\cref{sec:dpo} inverts: solving \cref{eq:rlhf-optimum} for
$r_\phi(x, y)$ in terms of $\pi^*$ and $\pi^{\mathrm{SFT}}$ gives
$r(x, y) = \beta \log \pi^*(y \mid x) / \pi^{\mathrm{SFT}}(y \mid x) +
\beta \log Z(x)$, and substituting this into the Bradley--Terry
preference probability \cref{eq:rlhf-rm} cancels the intractable
$\log Z(x)$, yielding the \glsterm{dpo}{DPO} loss as a closed-form classification
objective on the same preference data the RM in \cref{sec:rlhf-rm}
was trained on. The result: same fixed point, no sampling, no reward
model, no \glsterm{value}{value} model, no RL loop. \cref{sec:dpo} works through the
derivation and the variant family (\glsterm{ipo}{IPO}, \glsterm{kto}{KTO}, \glsterm{orpo}{ORPO}, \glsterm{simpo}{SimPO}, CPO).
